\documentclass[11pt,a4paper]{article}

% --- 1. Configuración de Página y Márgenes ---
\usepackage[margin=1in]{geometry}
\usepackage{authblk} % Gestión elegante de autores y afiliaciones

% --- 2. Paquetes Matemáticos ---
\usepackage{amsmath, amssymb, amsthm}
\usepackage{mathtools}
\usepackage{bm}
\usepackage{dsfont}
\usepackage{physics}
\usepackage{braket}

% --- 3. Elementos Gráficos y Tablas ---
\usepackage{graphicx}
\usepackage{color}
\usepackage{xcolor}
\usepackage{booktabs} % Líneas horizontales de alta calidad para tablas

% --- 4. Algoritmos y Código ---
\usepackage{algorithm}
\usepackage{algpseudocode}
\usepackage{listings}

% --- 5. Hipervínculos e Invariantes PDF ---
\usepackage[numbers,sort&compress]{natbib} 
\usepackage[colorlinks=true, linkcolor=blue!70!black, citecolor=red!70!black, urlcolor=blue!70!black]{hyperref}
\usepackage{orcidlink}
\newcommand{\orcid}[1]{\orcidlink{#1}}
\usepackage{url}

% --- 6. Cleveref (SIEMPRE el último del preámbulo) ---
\usepackage{cleveref}
\crefname{equation}{Eq.}{Eqs.}
\crefname{section}{Sec.}{Secs.}
\crefname{subsection}{Sec.}{Secs.}
\crefname{table}{Table}{Tables}
\crefname{figure}{Fig.}{Figs.}
\labelformat{lemma}{Lemma~#1}
\labelformat{theorem}{Theorem~#1}

% --- Configuración de Estilos de Código ---
\definecolor{codegreen}{rgb}{0,0.6,0}
\definecolor{codegray}{rgb}{0.5,0.5,0.5}
\definecolor{codepurple}{rgb}{0.58,0,0.82}
\definecolor{backcolour}{rgb}{0.95,0.95,0.92}

\lstdefinestyle{mystyle}{
    backgroundcolor=\color{backcolour},   
    commentstyle=\color{codegreen},
    keywordstyle=\color{magenta},
    numberstyle=\tiny\color{codegray},
    stringstyle=\color{codepurple},
    basicstyle=\ttfamily\footnotesize,
    breakatwhitespace=false,         
    breaklines=true,                  
    captionpos=b,                    
    keepspaces=true,                  
    numbers=left,                    
    numbersep=5pt,                  
    showspaces=false,                
    showstringspaces=false,
    showtabs=false,                  
    tabsize=2,
    language=Python
}
\lstset{style=mystyle}

% --- Entornos de teoremas ---
\newtheorem{theorem}{Theorem}[section]
\newtheorem{lemma}[theorem]{Lemma}
\newtheorem{corollary}[theorem]{Corollary}
\newtheorem{definition}[theorem]{Definition}
\newtheorem{remark}[theorem]{Remark}

% --- Macros personalizados ---
\newcommand{\Z}{\mathbb{Z}}
\newcommand{\C}{\mathbb{C}}
\newcommand{\R}{\mathbb{R}}
\newcommand{\N}{\mathbb{N}}
\newcommand{\Hilb}{\mathcal{H}}
\newcommand{\Zmod}{\mathbb{Z}/6\mathbb{Z}}
\DeclareMathOperator{\IPR}{IPR}

\begin{document}

\title{\textbf{Multifractal non-ergodic extended phase in power-law random banded matrices with modular arithmetic constraints}}

\author{José Ignacio Peinador Sala \orcidlink{0009-0008-1822-3452}}
\affil{Independent Researcher, Valladolid, Spain\\ \texttt{joseignacio.peinador@gmail.com}}

\date{\today}

\maketitle

\begin{abstract}
We introduce a hybrid deterministic-stochastic Hamiltonian for a one-dimensional quantum lattice, where off-diagonal hoppings are constrained by coprimality with respect to~6 (encoded in the ring $\mathbb{Z}/6\mathbb{Z}$). The diagonal potential implements the exact inversion of the Riemann--von Mangoldt counting function via the Lambert $W$ function, generating a smooth macroscopic density without adjustable parameters. The off-diagonal disorder belongs to the class of power-law random banded matrices (PRBM) with a decay exponent $\nu=0.75$, positioning the system at the center of the extended chaotic phase and ensuring essential self-adjointness. The chaos coupling is analytically set to $\epsilon = \pi\sqrt{2}$ via a local resonance condition. Exact diagonalization up to $N=20\,000$ demonstrates that the macroscopic spectral density matches the input Weyl law ($R^2=0.999997$), while local fluctuations follow Wigner--Dyson statistics of the Gaussian unitary ensemble (GUE). However, the long-range spectral rigidity, analyzed via an ensemble-averaged spectral form factor ($N=15\,000$), reveals a sub-diffusive ramp with exponent $\gamma = 0.6086 \pm 0.0103$. The corresponding generalized fractal dimension is $D_2 = 0.24338 \pm 0.00006$, confirming that the eigenstates occupy a sparse multifractal support. Numerical analysis shows that both $D_2$ and $\gamma$ scale with the density of surviving hopping channels $\varphi(m)/m$, placing the system in a non-ergodic extended (NEE) phase. This arithmetic mask provides topological protection against Anderson localization, preserving spectral rigidity beyond the unconstrained critical point. Furthermore, we prove a rigorous upper bound $D_2 \le \varphi(m)/m$, formally certified using the Lean~4 theorem prover, linking the algorithmic complexity of the prime sieve to the multifractal spatial support. A single-shot GPU diagonalization at $N=16\,000$ yields a bulk-averaged $D_2 = 0.2465$, lying just $1.4\%$ below the certified bipartite limit ($D_2 = 1/4$). Finally, we compute the spectral form factor of the first $10^4$ Riemann zeros, finding $\gamma \approx 1.12$ (consistent with GUE ergodicity), demonstrating that the observed multifractality is intrinsic to the modular constraint rather than the zeta spectrum. This model provides an analytically tractable example of an arithmetic NEE phase, with broader implications for noncommutative geometry and the engineering of robust quantum information channels.

\vspace{1em}
\noindent \textbf{Keywords:} Quantum Chaos, Non-ergodic Extended Phase, Power-law Random Banded Matrices, Noncommutative Geometry, Anderson Localization, Formal Verification.
\end{abstract}

\section{Introduction}
\label{sec:introduction}

Random matrix theory (RMT) and the physics of disordered quantum systems have long been intertwined with number theory, most famously through the statistical properties of the Riemann zeta zeros. Montgomery's pair correlation conjecture and Odlyzko's numerical computations demonstrated that the zeros on the critical line exhibit the same level repulsion as eigenvalues of Gaussian unitary ensemble (GUE) matrices \cite{Montgomery1973,Odlyzko2001}. This motivated the Hilbert--P\'olya idea that the imaginary parts of the zeros might be eigenvalues of a Hermitian operator, a program that remains open.

Rather than attempting to construct such an operator, we take a different route: we use the smooth density of zeros---the Riemann--von Mangoldt formula---as the diagonal potential of a one-dimensional lattice Hamiltonian, and then add off-diagonal disorder subject to an arithmetic constraint derived from the ring $\mathbb{Z}/6\mathbb{Z}$. The resulting model is a deterministic-stochastic hybrid whose thermodynamic and dynamical properties can be studied in their own right.

The primary structural feature of this model is the arithmetic selection rule: hopping between sites $i$ and $j$ is allowed only if the distance $|i-j|$ is coprime to~6, meaning it belongs to the residue classes $1$ or $5$ modulo~6. This rule eliminates exactly two-thirds of the connections and introduces a periodic sparsification of the interaction graph. The system falls into the universality class of power-law random banded matrices (PRBM) with spatial decay $\nu=0.75$, a value that places it at the center of the extended chaotic phase where eigenstates are neither fully ergodic nor localized \cite{Mirlin1996}.

We show by large-scale exact diagonalization that the macroscopic spectral density remains locked to the input Weyl law, while the local level statistics are GUE. However, the long-range spectral rigidity, quantified by the spectral form factor (SFF) and the generalized fractal dimension $D_2$, deviates systematically from the ergodic limit. The SFF ramp exponent is $\gamma \approx 0.609$, below the standard GUE value $\gamma=1$, and the fractal dimension is $D_2 \approx 0.243$, indicating that wavefunctions populate a sparse multifractal support. These independent metrics consistently place the system in a non-ergodic extended (NEE) phase.

Understanding the mechanisms that generate NEE phases is of significant interest for the development of quantum information technologies, as these phases offer a balance between state extendedness and limited thermalization. We demonstrate that $D_2$ and $\gamma$ are systematically suppressed as the density of surviving hopping channels $\varphi(m)/m$ decreases. This channel-density scaling, together with a rigorous upper bound $D_2 \le \varphi(m)/m$, establishes a direct connection between the arithmetic restriction and the emergent multifractality. Furthermore, we show that the arithmetic mask provides topological protection against Anderson localization, preserving GUE spectral rigidity well beyond the unconstrained critical point. This suggests that modular arithmetic filters could be utilized to engineer robust quantum networks that resist disorder-induced localization.

To test whether this sub-diffusive behavior might reflect an unknown property of the Riemann zeros themselves, we compute the SFF of the first $10^4$ actual zeros. We find $\gamma \approx 1.12$, which is consistent with ordinary GUE. This discrepancy confirms that the observed NEE phase is an intrinsic feature of the modular constraint and the PRBM structure, not a spectral fingerprint of the zeta function.

The model presented here serves as a mathematically tractable example of how a simple arithmetic filter generates multifractal non-ergodic behavior in a quantum system. The $\mathbb{Z}/6\mathbb{Z}$ mask is not an arbitrary choice; it emerges naturally as the center of the Standard Model gauge group in noncommutative geometry, where the internal space has KO-dimension 6 modulo 8 \cite{Connes2008,Chamseddine2007}. This geometric origin provides a robust theoretical motivation for the arithmetic constraint.

The paper is organized as follows. \Cref{sec:model} defines the Hamiltonian and its parameters. \Cref{sec:results} presents the numerical results, including the macroscopic alignment, the GUE level statistics, the SFF and $D_2$ analysis, the channel-density scaling, the topological protection, and the comparison with the actual Riemann zeros. \Cref{sec:discussion} details the physical interpretation and the connection to noncommutative geometry. Technical derivations are collected in the Appendix, including the formal bound $D_2 \le \varphi(m)/m$.

\section{Model Hamiltonian}
\label{sec:model}

\subsection{Hilbert space and diagonal potential}

We consider a quantum system on the Hilbert space $\Hilb = \ell^2(\mathbb{N})$ with orthonormal basis $\{|k\rangle\}_{k\ge 1}$. The unperturbed Hamiltonian is diagonal,
\begin{equation}
\hat H_0 |k\rangle = E_k |k\rangle,
\end{equation}
with the energies $E_k$ chosen to reproduce the macroscopic density of a generic sequence with a logarithmic level density, such as the Riemann zeros. The Riemann--von Mangoldt formula states that the number of zeros with imaginary part less than $E$ is
\begin{equation}
N(E) = \frac{E}{2\pi}\log\left(\frac{E}{2\pi e}\right) + \frac{7}{8} + O(E^{-1}).
\end{equation}
Setting $N(E_k)=k$ and inverting exactly by the Lambert $W$ function gives the deterministic diagonal term
\begin{equation}
\boxed{ E_k = \frac{2\pi (k - 7/8)}{W\!\left(\frac{k-7/8}{e}\right)} }.
\label{eq:diag}
\end{equation}
The constant $7/8$ is the Maslov phase originating from the Riemann--Siegel theta function \cite{Corless1996}. In the unfolded spectrum, the mean level spacing tends to $2\pi$.

\subsection{Off-diagonal disorder and modular constraint}

The perturbation $\hat V$ is a random matrix with elements
\begin{equation}
V_{ij} = \epsilon \, d^{-\nu} \, G_{ij} \, \Xi(d),
\qquad d = |i-j|,
\end{equation}
where $G_{ij}$ are independent complex Gaussian variables with $\mathbb{E}[|G_{ij}|^2]=1$, $\epsilon$ is the coupling constant, $\nu$ the spatial decay exponent, and $\Xi(d)$ is the arithmetic mask
\begin{equation}
\Xi(d) = \begin{cases}
1 & \text{if } \gcd(d,6)=1,\\
0 & \text{otherwise}.
\end{cases}
\end{equation}
This mask restricts hopping to distances $d\equiv1,5\pmod6$, thereby blocking all connections whose length shares a factor with~6. The total Hamiltonian
\begin{equation}
\hat H = \hat H_0 + \hat V
\label{eq:fullH}
\end{equation}
is manifestly Hermitian because $V_{ji} = \overline{V_{ij}}$.

\subsection{Parameters: \texorpdfstring{$\nu$}{nu}, \texorpdfstring{$\epsilon$}{epsilon}, and self-adjointness}

The operator $\hat V$ belongs to the PRBM class. The Kato--Rellich theorem requires that $\hat V$ be relatively bounded with respect to $\hat H_0$. Since $\sum_d d^{-2\nu}$ converges only for $\nu>0.5$, the condition $\nu>0.5$ must be met. To remain within the extended chaotic phase and avoid both thermal divergence ($\nu\le0.5$) and Anderson localization ($\nu\ge1$), we fix
\begin{equation}
\nu = 0.75,
\end{equation}
the exact center of the critical window $0.5<\nu<1$ \cite{Mirlin1996}. This value guarantees essential self-adjointness on the dense domain $\mathcal{D}(\hat H_0)$.

The coupling $\epsilon$ is determined by a local resonance condition. In the unfolded spectrum, the mean spacing is $D=2\pi$. Two neighboring levels hybridize effectively when the variance of the coupling compensates the level spacing. Because the interaction is Hermitian, a transition $|n\rangle\to|m\rangle$ has a conjugate return path, doubling the effective variance in the local $2\times2$ subspace. Hence we impose
\begin{equation}
2\,\mathbb{E}[|V_{n,n+1}|^2] = D^2 \;\Longrightarrow\; 2\epsilon^2 = (2\pi)^2,
\end{equation}
yielding
\begin{equation}
\boxed{\epsilon_c = \pi\sqrt{2} \approx 4.44288}.
\label{eq:epsilon}
\end{equation}
This value lies at the threshold of maximal chaos in the local subspace \cite{Cunden2018,Gross1980}.

\subsection{Connection to noncommutative geometry}

While the model is mathematically self-contained, the choice of modulus~6 is grounded in physical theory. In Connes' noncommutative geometry approach to the Standard Model, the finite internal space is described by the algebra $\mathcal{A}_F = \mathbb{C}\oplus\mathbb{H}\oplus M_3(\mathbb{C})$, whose unimodular gauge group is $G_{SM} = U(1)\times SU(2)\times SU(3)/\mu_6$. The center of this group is exactly $\mathbb{Z}/6\mathbb{Z}$ \cite{Connes2008}. To avoid fermion doubling and produce the correct chiral spectrum, the internal space must have KO-dimension $6$ modulo $8$ \cite{Chamseddine2007}. The arithmetic mask, which selects the two residues coprime to~6, functions as a discrete realization of the chiral grading of the internal space. 

\subsection{Discrete Dirac operator and the \texorpdfstring{$\mathbb{Z}/6\mathbb{Z}$}{Z/6Z} superselection rule}
\label{sec:ncg_discrete}

The connection between the modular ring $\Zmod$ and the finite geometry of the Standard Model can be made quantitatively precise by considering the discretization of the Dirac operator on a one‑dimensional lattice that inherits the algebraic structure of the internal space $F$. 

In the spectral triple $(\mathcal{A}, \mathcal{H}, D)$ describing the full Standard Model \cite{Connes2008,Chamseddine2007}, the algebra is $\mathcal{A}=C^\infty(M)\otimes \mathcal{A}_F$ with $\mathcal{A}_F=\mathbb{C}\oplus\mathbb{H}\oplus M_3(\mathbb{C})$, and the Dirac operator decomposes as $D = D_M\otimes 1 + \gamma_M\otimes D_F$. The finite part $D_F$ acts on the fermionic Hilbert space $\mathcal{H}_F$ and encodes the Yukawa couplings and the CKM matrix.

The unimodular gauge group is $G_{SM}=U(1)\times SU(2)\times SU(3)\,/\,\mu_6$, where the discrete group $\mu_6\simeq\Zmod$ acts by multiplication with sixth roots of unity. This quotient is necessary to cancel the common center of the three factors and to obtain the correct hypercharge assignments. Physical fields must be invariant under this central action, imposing a superselection rule: any transition between two states whose $U(1)$ charges differ by an integer that is not a multiple of~6 is forbidden.

Translating this restriction into a one‑dimensional lattice model involves replacing the continuous spacetime $M$ by a discrete circle $\mathbb{Z}_L$ with $L\to\infty$, and projecting the internal space onto a single direction carrying charge information. Let $\psi_n$ be a fermionic amplitude localized at site $n\in\mathbb{Z}$. The Dirac operator on the lattice is written as
\begin{equation}
(D\psi)_n = \sum_{d\neq 0} \frac{\gamma(d)}{d^{\alpha}} \, \psi_{n+d} + \text{diagonal terms},
\end{equation}
where $\gamma(d)$ are matrices in the internal space and the power‑law decay $d^{-\alpha}$ mimics the long‑range nature of the Euclidean propagator on a one‑dimensional manifold (see Ref.~\cite{Mirlin1996}).

Invariance under the central $\mu_6$ forces the transition amplitude between sites $n$ and $n+d$ to carry a phase factor that cancels the action of $\mu_6$ on the internal indices. In the basis where the $U(1)$ generator is diagonal, the condition is
\begin{equation}
e^{2\pi i q d /6} = 1 \quad \text{for all allowed transitions},
\end{equation}
where $q$ is the $U(1)$ charge of the fermion. This equation requires $\gcd(d,6)=1$, allowing the phase to be absorbed by a suitable choice of the internal matrix $\gamma(d)$. Hence, off‑diagonal entries are restricted to distances $d$ coprime to~6. The general charge‑invariant kinetic term is
\begin{equation}
(D\psi)_n = \sum_{\substack{d\neq 0 \\ \gcd(d,6)=1}} \frac{\Gamma_d}{d^{\nu}} \, \psi_{n+d} + \text{h.c.},
\end{equation}
where $\Gamma_d$ are complex random matrices (GUE) encoding the remaining geometric information of the finite space $F$. Identifying the coupling $\epsilon$ with the overall scale of $D_F$ and $\nu$ with the fall‑off exponent of the propagator recovers the off‑diagonal part of Eq.~\eqref{eq:fullH}.

The diagonal potential $E_k$ in Eq.~\eqref{eq:diag} is similarly motivated from the NCG framework. The finite Dirac operator $D_F$ has a block‑diagonal form whose eigenvalues are related to fermion masses. After integrating out continuous degrees of freedom and taking the logarithmic large‑$N$ limit, the effective eigenvalue density becomes proportional to the Riemann‑von Mangoldt formula \cite{Connes1999}. The exact Lambert $W$ inversion used here is a discrete analogue of this asymptotic density, with the $7/8$ shift acting as the Maslov index that corrects boundary terms in the trace formula.

Therefore, the full Hamiltonian $\hat H = \hat H_0 + \hat V$ can be regarded as a finite‑dimensional approximation of the Dirac operator $D$ projected onto a one‑dimensional lattice that retains the essential $\mathbb{Z}/6\mathbb{Z}$ central invariance. The arithmetic mask $\Xi(d)=\chi_0^{(6)}(d)$ serves as the direct imprint of the KO-dimension 6 chiral grading on the discrete kinetic term. 

\section{Numerical results}
\label{sec:results}

\subsection{Exact diagonalization and macroscopic alignment}

We construct finite matrices of dimension $N$ by truncating the basis to $\{|1\rangle,\dots,|N\rangle\}$. The diagonal entries are given by Eq.~\eqref{eq:diag} with $k = 2,\dots,N+1$, and the off-diagonal entries are populated as described above. The matrix is dense in the allowed channels but structurally sparse: only one-third of the distances survive the mask. Diagonalization is performed with \texttt{scipy.linalg.eigh} using complex64 precision. For the main runs, we evaluate $N=20\,000$ and compare the first $10\,000$ eigenvalues with the corresponding Riemann zeros (used exclusively as a benchmark for the macroscopic density).

Figure~\ref{fig:macroscopic} (left panel) shows the raw eigenvalues alongside the Riemann zeros. The coefficient of determination is $R^2 = 0.999997$, and the macroscopic scale factor converges autonomously to $0.9994$, within $0.06\%$ of unity. This confirms that the Lambert $W$ initialization accurately implements the desired Weyl law without empirical rescaling.

\begin{figure}[htbp]
\centering
\includegraphics[width=0.8\textwidth]{PRL_Figure_Ultimate_10k.png}
\caption{\textbf{Left:} Eigenvalues of $\hat H$ (blue dashed) vs.\ first $10^4$ Riemann zeros (red). \textbf{Center:} Scatter plot showing $R^2 = 0.999997$. \textbf{Right:} Unfolded nearest-neighbor spacing distribution (green) vs.\ GUE Wigner surmise (black).}
\label{fig:macroscopic}
\end{figure}

\subsection{Local spectral statistics}

We unfold the spectrum using the theoretical counting function
\begin{equation}
w_n = \frac{E_n}{2\pi}\log\left(\frac{E_n}{2\pi e}\right),
\end{equation}
and compute the nearest-neighbor spacing distribution $P(s)$. Figure~\ref{fig:macroscopic} (right panel) demonstrates alignment with the GUE Wigner surmise $P(s) = \frac{32}{\pi^2}s^2 e^{-4s^2/\pi}$. The level repulsion $P(0)\to0$ is evident, excluding Poissonian integrable statistics.

\subsection{Spectral form factor and \texorpdfstring{$D_2$}{D2}: evidence for the NEE phase}

The spectral form factor (SFF) probes long-range correlations and is defined as
\begin{equation}
K(t) = \frac{1}{N}\Bigl|\sum_{n=1}^N e^{-i w_n t}\Bigr|^2.
\end{equation}
Because the SFF is not self-averaging, we perform an ensemble average over $M=100$ independent realizations of the GUE noise, each of size $N=15\,000$. The averaged SFF $\langle K(t)\rangle$ is displayed in Fig.~\ref{fig:sff}. It exhibits the canonical ``dip--ramp--plateau'' structure, but the ramp (region $0.5 \lesssim t \lesssim 4$) follows a power law $\langle K(t)\rangle \propto t^\gamma$ with
\begin{equation}
\gamma = 0.6086 \pm 0.0103 \quad (95\%\ \text{CI } [0.5835,0.6328]).
\end{equation}
This value is significantly smaller than the ergodic GUE prediction of $\gamma=1$. The plateau saturates at $\langle K(t)\rangle =1$ for $t>2\pi$, confirming the Hermitian nature of the spectrum.

For the same eigenstates, we compute the inverse participation ratio $\IPR_n = \sum_i |\psi_n(i)|^4$ and extract the generalized fractal dimension $D_2$ from the scaling $\IPR \sim N^{-D_2}$. Averaging over the central 50\% of the spectrum and over the $M=100$ realizations yields
\begin{equation}
D_2 = 0.24338 \pm 0.00006.
\end{equation}
A value of $D_2 < 1$ indicates that wavefunctions do not uniformly explore the full Hilbert space but reside on a fractal support of reduced dimension. 

Together, $\gamma<1$ and $D_2<1$ define the Non-Ergodic Extended (NEE) phase: the states are extended (the system is not localized, as the level statistics are GUE) but do not thermalize over the available space. 

\begin{figure}[htbp]
\centering
\includegraphics[width=0.8\textwidth]{PRL_Figure_Final_con_inset.png}
\caption{Ensemble-averaged spectral form factor $\langle K(t)\rangle$ ($M=100$, $N=15\,000$). The ramp (blue solid) follows a power law with exponent $\gamma \approx 0.609$, compared to the GUE prediction $\gamma=1$ (dashed). Inset: log-log magnification of the ramp region.}
\label{fig:sff}
\end{figure}

\subsection{Channel-density scaling of \texorpdfstring{$D_2$}{D2} and \texorpdfstring{$\gamma$}{gamma}}
\label{sec:channel_density_scaling}

To isolate the effect of the arithmetic restriction, we compare four versions of the Hamiltonian with different modular masks $m$, corresponding to allowed distances satisfying $\gcd(d,m)=1$. The fraction of surviving channels is $\varphi(m)/m$, where $\varphi$ is Euler's totient function. The results for $N=2\,000$ and $\nu=0.75$ (averaged over $M=5$ realizations) are detailed in Table~\ref{tab:mask_comparison}.

\begin{table}[htbp]
\centering
\caption{Systematic suppression of $D_2$ and $\gamma$ with increasing arithmetic restriction ($N=2\,000$, $\nu=0.75$, $M=5$).}
\label{tab:mask_comparison}
\begin{tabular}{c c c c c}
\toprule
\textbf{Mask} & $\varphi(m)/m$ & $D_2$ & $\gamma$ & \textbf{Interpretation} \\
\midrule
None (PRBM)     & $1.0000$ & $0.3438 \pm 0.0013$ & $0.7191$ & Mildly restricted chaos \\
$\mathbb{Z}/2\mathbb{Z}$ (odd $d$) & $0.5000$ & $0.3078 \pm 0.0016$ & $0.7440$ & Parity alone insufficient \\
$\mathbb{Z}/6\mathbb{Z}$ (this work) & $0.3333$ & $0.2796 \pm 0.0017$ & $0.4816$ & Strong NEE signature \\
$\mathbb{Z}/30\mathbb{Z}$ & $0.2667$ & $0.2551 \pm 0.0022$ & $0.4013$ & Deeper into NEE phase \\
\bottomrule
\end{tabular}
\end{table}

First, $D_2$ decreases monotonically from $0.344$ (no mask) to $0.255$ ($m=30$), tracking $\varphi(m)/m$. This indicates that the spatial support of the wavefunctions is controlled by the density of allowed hopping channels, independently of the PRBM decay exponent $\nu$. Second, $\gamma$ exhibits a non-monotonic step at $m=2$ ($0.719 \to 0.744$) before decreasing sharply to $0.482$ for $m=6$ and $0.401$ for $m=30$. The parity mask alone is insufficient to break ergodicity; the elimination of both prime factors 2 and 3 is the minimal condition to trigger sub-diffusive dynamics. 

Finite-size scaling analysis shows that $D_2$ and $\gamma$ converge to stable asymptotic values in the thermodynamic limit. Table~\ref{tab:scaling} tracks the evolution with system size for the $\mathbb{Z}/6\mathbb{Z}$ mask, using $M=5$ realizations for $N \le 8\,000$ and $M=100$ for $N=15\,000$. $D_2$ decreases from $0.278$ to $0.243$, bounded above by $\varphi(6)/6 = 1/3$, while $\gamma$ increases from $0.531$ to $0.610$, saturating below unity. These results support the thermodynamic stability of the NEE phase.

\begin{table}[htbp]
\centering
\caption{Finite-size scaling of $D_2$ and $\gamma$ for the $\mathbb{Z}/6\mathbb{Z}$ mask ($\nu=0.75$).}
\label{tab:scaling}
\begin{tabular}{c c c c}
\toprule
$N$ & $D_2$ & $\gamma$ & $M$ \\
\midrule
$2\,000$  & $0.2781 \pm 0.0027$ & $0.5314$ & $5$ \\
$4\,000$  & $0.2650 \pm 0.0010$ & $0.5398$ & $5$ \\
$8\,000$  & $0.2532 \pm 0.0012$ & $0.5758$ & $5$ \\
$15\,000$ & $0.2433 \pm 0.0006$ & $0.6100 \pm 0.0099$ & $100$ \\
\bottomrule
\end{tabular}
\end{table}

\subsection{Robustness of the chaotic phase}

The chaotic phase supported by the arithmetic mask exhibits structural stability. Figure~\ref{fig:robustness} shows the average level spacing ratio $\langle r \rangle$ as a function of the coupling strength $\epsilon$ for $N=1\,000$ and $\nu=0.75$. The system enters the GUE regime ($\langle r \rangle \approx 0.58$--$0.60$) at $\epsilon \approx 3$ and remains chaotic throughout the upper range. The theoretical value $\epsilon_c = \pi\sqrt{2} \approx 4.44$ lies within this plateau, establishing that the GUE universality class is not fine-tuned.

\begin{figure}[htbp]
\centering
\includegraphics[width=1\textwidth]{Robustness_Epsilon_Nu.png}
\caption{Robustness of the chaotic phase under the $\mathbb{Z}/6\mathbb{Z}$ mask ($N=1\,000$). \textbf{Left:} Level spacing ratio $\langle r \rangle$ vs.\ coupling strength $\epsilon$ ($\nu=0.75$). The system exhibits a broad chaotic plateau; the vertical dashed line marks the theoretical resonance $\epsilon_c = \pi\sqrt{2}$. \textbf{Right:} Level spacing ratio $\langle r \rangle$ vs.\ spatial decay $\nu$ ($\epsilon = \pi\sqrt{2}$). The arithmetic mask provides topological protection against Anderson localization, maintaining $\langle r \rangle > 0.53$ at $\nu=1.2$, above the Poisson limit ($0.386$, dashed red) to which the unmasked standard PRBM (dashed red line) decays.}
\label{fig:robustness}
\end{figure}

The arithmetic mask provides topological protection against Anderson localization. In standard PRBM, a transition from Wigner--Dyson ($\langle r \rangle \approx 0.60$) to Poisson ($\langle r \rangle \approx 0.386$) statistics occurs at $\nu_c = 1.0$ \cite{Mirlin1996}. Computing $\langle r \rangle$ for the model with the $\mathbb{Z}/6\mathbb{Z}$ mask over $\nu \in [0.5, 1.2]$ at $N=1\,500$ (Fig.~\ref{fig:robustness}), $\langle r \rangle$ decays from $0.60$ at $\nu=0.5$ to $0.54$ at $\nu=1.2$. This persistence of spectral rigidity beyond the unconstrained critical point indicates that the arithmetic mask establishes long-range entangled channels that sustain quantum communication even when the bare kinetic coupling would localize.

\subsection{Microscopic multifractal spectrum: GPU diagonalization}

To verify that the multifractal support is a microscopic property of individual eigenstates rather than an artifact of ensemble averaging, we performed a single-shot exact diagonalization of $\hat H$ at $N=16\,000$ on a GPU (NVIDIA T4). The generalized fractal dimension $D_2^{(n)}$ was computed for each eigenstate $n$. Figure~\ref{fig:fractal} displays $D_2^{(n)}$ for the full spectrum and a histogram of the central 60\% bulk states.

The bulk-averaged dimension is
\begin{equation}
\langle D_2 \rangle_{\text{bulk}} = 0.2465 \pm 0.0313,
\end{equation}
with a median of $0.2505$ and an interquartile range of $[0.2280, 0.2694]$. These values align with the theoretical limit $D_2 = 1/4$ derived from the product-of-memories bound (\Cref{app:bound}), with a relative error of $1.4\%$. The distribution is localized around $0.25$ and remains separate from the ergodic reference $D_2=1$. This confirms that the arithmetic mask confines the eigenstates to a low-dimensional fractal subspace set by the modular totient fraction.

\begin{figure}[htbp]
\centering
\includegraphics[width=0.8\textwidth]{Fractal_Dimension_D2.png}
\caption{Microscopic multifractal support of $\hat H$ at $N=16\,000$. Left: $D_2^{(n)}$ for every eigenstate $n$ (teal points). The horizontal lines mark the ergodic limit ($D_2=1$, red dashed), the theoretical bound $D_2=0.25$ (orange solid), and the bulk average (blue dotted). Right: probability density of $D_2$ in the central 60\% of the spectrum. The distribution is localized around $0.25$, consistent with the NEE phase for individual eigenstates.}
\label{fig:fractal}
\end{figure}

\subsection{Formal mechanized certification of arithmetic bounds}
\label{sec:formal_verification}

While exact diagonalization reveals the physical manifestations of the non-ergodic extended phase (the sub-diffusive ramp $\gamma$ and the fractal dimension $D_2$), the fundamental arithmetic rules governing the sparsity of the interaction graph can be proven axiomatically. To ensure mathematical rigor regarding the asymptotic limits, we formalized the core discrete topological constraints using the Lean~4 theorem prover.

This mechanized verification certifies four properties governing the NEE phase. First, we verified the critical exponent arithmetic, confirming that for $\nu=0.75$, the unweighted degree scaling exponent $1-2\nu$ evaluates to $-0.5$, satisfying the convergence requirements of the Kato-Rellich theorem for essential self-adjointness. Second, we formally proved the totient channel density bound, showing that the maximum generic fractional dimension $D_2 \le \varphi(m)/m$ evaluates exactly to $1/3$ for the modular mask $\mathbb{Z}/6\mathbb{Z}$.

Third, we certified the graph sparsity base block, proving that the arithmetic mask allows exactly 2 valid transitions in every sequence of 6 lattice sites, forming the discrete topological foundation for the expected degree asymptotic limit utilized in our analytical derivations (see \Cref{app:bound}). Finally, we verified the algebraic consistency of the bipartite factorization bound; if the topological space divides symmetrically into two chiral sublattices, the maximum factorized probability density evaluates exactly to $0.25$.

The bulk-averaged empirical dimension from the GPU diagonalization ($\langle D_2 \rangle_{\text{bulk}} = 0.2465$) lies within $1.4\%$ of this formally certified limit ($1/4$), providing strong evidence that the multifractal geometry of the wavefunctions is dictated by these arithmetic axioms. The Lean~4 environment and proof scripts are provided as a supplementary interactive notebook.

\subsection{Comparison with the Riemann zeros}

The diagonal potential Eq.~\eqref{eq:diag} was constructed using the Riemann zeros. To determine if the NEE phase is a property of the actual zeta zeros, we computed the SFF of the first $10\,000$ Riemann zeros directly, using the same unfolding and block-averaging procedure (10 blocks of 1\,000 levels each). The resulting ramp exponent is
\begin{equation}
\gamma_{\text{zeros}} = 1.1159 \pm 0.0294 \quad (R^2=0.78),
\end{equation}
which is consistent with $\gamma=1$ and the conventional GUE picture. The sub-diffusive behavior observed in the model is an intrinsic property of the modular PRBM Hamiltonian, not a feature of the true Riemann spectrum.

\section{Discussion}
\label{sec:discussion}

The results demonstrate that a minimal arithmetic constraint---coprimality with respect to~6---applied to the hoppings of a PRBM drives the system into a robust NEE phase. The phase exhibits GUE local statistics, a sub-diffusive SFF ramp ($\gamma\approx0.61$), and a multifractal eigenstate structure ($D_2\approx0.24$). These values are stable under finite-size scaling and across independent disorder realizations.

The parameters $\nu$ and $\epsilon$ are fixed by mathematical constraints rather than curve fitting. The decay exponent $\nu=0.75$ lies at the center of the PRBM extended phase, and the coupling $\epsilon=\pi\sqrt{2}$ arises from a local resonance condition analogous to the Thouless criterion. The mask $\Xi(d)$ originates from geometric principles: $\mathbb{Z}/6\mathbb{Z}$ is the center of the Standard Model gauge group in the noncommutative geometry framework, and the KO-dimension 6 condition selects the same modular structure. This allows the model to act as a discretized representation of a chiral gauge theory on a one-dimensional lattice.

The channel-density scaling of $D_2$ and $\gamma$ (\Cref{sec:channel_density_scaling}) demonstrates that the NEE phase is a systematic consequence of arithmetic sparsification. The monotonic decrease of $D_2$ with $\varphi(m)/m$ links the fractal support to the number of allowed hopping channels. The parity filter ($m=2$) does not induce the NEE phase, highlighting that the elimination of both prime factors 2 and 3 (the densest sources of composite noise) generates the necessary conditions for this phase.

The topological protection against Anderson localization (\Cref{fig:robustness}) and the robustness of chaos over a broad range of $\epsilon$ confirm that the chaotic phase is a structural feature of the arithmetic mask. The arithmetic constraint provides a backbone of long-range correlations that resists the localization transition. This finding holds practical implications for fields such as quantum information science and materials engineering; employing specific arithmetic or modular constraints could aid in the design of robust quantum communication networks capable of maintaining coherence in disordered environments.

A rigorous upper bound $D_2 \le \varphi(m)/m$ is derived in \Cref{app:bound} from the expected degree of the arithmetic interaction graph. This bound reflects an information-theoretic constraint: the algorithmic memory required to generate the prime sieve scales as $O(\sqrt{N})$, restricting the fractal support to a fraction of the full Hilbert space. For $m=6$, the bound gives $D_2 \le 1/3$, and the stronger product-of-memories argument refines this to $D_2 \le 1/4$. The GPU diagonalization at $N=16\,000$ yields a bulk-averaged $D_2 = 0.2465$, which closely aligns with this bound, indicating that the arithmetic mask governs the fractal geometry of individual eigenstates.

The observation that the Riemann zeros do not exhibit the sub-diffusive ramp suggests that the true zeta zeros behave as eigenvalues of a dense GUE matrix at the spectral scales probed by the SFF. 

The NEE phase identified in this model adds to existing research on multifractal non-ergodic states in Anderson transitions on random graphs \cite{Kravtsov2015}, Sachdev-Ye-Kitaev models \cite{Legramandi2024}, and the Rosenzweig-Porter ensemble. This model provides an exactly solvable limit where non-ergodicity is enforced by an algebraic selection rule rather than quenched random potentials. The arithmetic origin of the sparsification suggests potential avenues for studying spectral properties of $L$-functions over finite fields.

\section{Conclusions}
\label{sec:conclusions}

We have constructed a hybrid deterministic-stochastic Hamiltonian combining an exact inverse Weyl law potential with a modular arithmetic mask and PRBM long-range hopping. The system demonstrates GUE level repulsion alongside sub-diffusive spectral rigidity and strong wavefunction multifractality, placing it in the non-ergodic extended phase. The parameters are determined by physical and mathematical constraints. We have shown that $D_2$ and $\gamma$ scale with the density of allowed hopping channels, that the chaotic phase is robust across coupling strengths and resists Anderson localization, and that a rigorous upper bound $D_2 \le \varphi(m)/m$ connects the multifractality to the arithmetic constraint. The measurement of the actual Riemann zeros as ergodic clarifies that the modular constraint drives the NEE behavior in this model. A single-shot GPU diagonalization at $N=16\,000$ confirms that eigenstates individually respect the fractal bound $D_2\approx 1/4$. This research provides an analytically tractable model of arithmetic quantum chaos with a stable non-ergodic phase and establishes a framework for exploring modular lattices in quantum systems.

% ==============================================================================
% DECLARATIONS AND ACKNOWLEDGMENTS
% ==============================================================================

\section*{Acknowledgments}
The author conducted this research independently and acknowledges the open-source community for providing the computational tools necessary for this work. This research received no specific grant from any funding agency in the public, commercial, or not-for-profit sectors.

\section*{Data Availability Statement}
The data, Python source code, Jupyter notebooks for exact diagonalization, and the Lean~4 interactive environment for formal mechanized verification required to reproduce all experiments, proofs, and statistical analyses presented in this manuscript are openly available in the project's GitHub repository at \url{https://github.com/NachoPeinador/The-Riemann-GUE-Hamiltonian} and are permanently archived on Zenodo at \url{https://doi.org/10.5281/zenodo.19284510}.

\section*{Competing Interests}
The author declares no competing financial or non-financial interests.

\section*{Declaration of Generative AI and AI-Assisted Technologies}
During the preparation of this work, the author used generative AI (Google Gemini) strictly for English language translation, stylistic polishing of the manuscript, and syntax formatting of the Python code. After using this tool, the author reviewed and edited the content as needed and takes full responsibility for the theoretical framework, physical interpretations, and the overall integrity of the publication.

\appendix

\section{Derivation of \texorpdfstring{$\epsilon_c$}{epsilon\_c} and the fractional ramp}
\label{app:ramp}

\subsection{Resonance condition for \texorpdfstring{$\epsilon_c$}{epsilon\_c}}
The local level spacing in the unfolded spectrum is $D=2\pi$. For two neighboring levels to fully hybridize, the variance of the off-diagonal coupling must satisfy $2\sigma_V^2 = D^2$. With $\sigma_V^2 = \epsilon^2 \mathbb{E}[|G|^2] = \epsilon^2$, this yields $\epsilon_c = \pi\sqrt{2}$.

\subsection{Altshuler-Shklovskii scaling and the SFF ramp}
For a PRBM in the NEE phase, the mesoscopic variance of the level number scales as $\Sigma^2(\Delta E) \sim (\Delta E)^{D_2/d}$ for $d=1$. The SFF is the Fourier transform of the two-point correlation function, and therefore the connected ramp behaves as $K_{\text{conn}}(t) \propto t^{\gamma}$ with $\gamma = D_2 + \eta$, where $\eta>0$ accounts for enhanced backscattering due to the sparse connectivity. With $D_2\approx0.243$ and $\eta\approx0.365$, we obtain $\gamma\approx0.608$, in agreement with the numerical fit.

\section{Formal bound \texorpdfstring{$D_2 \le \varphi(m)/m$}{D2 <= phi(m)/m} from the arithmetic interaction graph}
\label{app:bound}

Let $\mathcal{G}_N$ be the graph whose vertices are the $N$ basis states and whose edges connect sites $i,j$ whenever $\gcd(|i-j|,m)=1$. The Hamiltonian is $\hat H = \hat H_0 + \hat V$, where $\hat V_{ij} = \epsilon_c \, d^{-\nu} \, G_{ij}$ on the edges and zero otherwise, with $\nu=0.75$. We define the generalized fractal dimension $D_2$ from the inverse participation ratio of the bulk eigenvectors, $D_2 = -\langle\log \IPR\rangle / \log N$, where the average is taken over the central half of the spectrum and over disorder realizations.

The foundational discrete arithmetic underlying the following graph limits, including the exact evaluation of the base block sparsity and the bipartite factorization, has been formally certified using the Lean~4 theorem prover (see \Cref{sec:formal_verification}).

\begin{lemma}[Expected degree]\label{lem:degree}
The expected number of non-zero off-diagonal entries in a row of $\hat V$ scales as
\begin{equation}
\langle \deg \rangle_N = \frac{\varphi(m)}{m} \, \frac{N^{1-2\nu}}{1-2\nu} + O(1).
\end{equation}
\end{lemma}
\begin{proof}
For a fixed site $i$, the number of sites $j$ with $|i-j|=d$ and $\gcd(d,m)=1$ is $2$ for each $d\le N-1$. Summing with the weight $d^{-2\nu}$ gives
\begin{align*}
\langle \deg \rangle_N &= 2\sum_{\substack{1\le d\le N-1 \\ \gcd(d,m)=1}} d^{-2\nu} \\
&= \frac{\varphi(m)}{m} \sum_{d=1}^{N-1} d^{-2\nu} + O(1) \\
&= \frac{\varphi(m)}{m} \frac{N^{1-2\nu}}{1-2\nu} + O(1),
\end{align*}
where we used that the fraction of integers coprime to $m$ is asymptotically $\varphi(m)/m$.
\end{proof}

\begin{lemma}[PRBM fractal dimension]\label{lem:prbm}
For a PRBM with exponent $\nu$ on a graph with average degree $k_N$, the generalized fractal dimension satisfies
\begin{equation}
D_2 = \min\!\left(1,\; \frac{\log k_N}{\log N}\right) + o(1)
\end{equation}
as $N\to\infty$, provided $1/2<\nu<1$.
\end{lemma}
\begin{proof}[Sketch]
In the NEE regime of PRBM, wavefunctions are multifractal with a support set of size $N^{D_2}$. The effective number of independent channels propagating quantum information is proportional to the conductance, which for $\nu<1$ is controlled by the average degree $k_N$ \cite{Mirlin1996}. Standard scaling theory gives $D_2 \sim \log(\text{conductance})/\log N \sim \log k_N / \log N$.
\end{proof}

\begin{theorem}[Channel-density bound]\label{thm:bound}
For the arithmetic Hamiltonian $\hat H$ with mask modulus $m$ and $\nu=0.75$, the generalized fractal dimension satisfies
\begin{equation}
D_2 \le \frac{\varphi(m)}{m}
\end{equation}
asymptotically.
\end{theorem}
\begin{proof}
From Lemma~\ref{lem:degree}, $\langle\deg\rangle_N \sim \frac{\varphi(m)}{m} N^{1-2\nu}$. With $\nu=0.75$, $1-2\nu = -0.5$, so $k_N \sim \frac{\varphi(m)}{m} N^{-0.5}$. Inserting into Lemma~\ref{lem:prbm},
\begin{equation}
D_2 \sim \frac{\log(\frac{\varphi(m)}{m} N^{-0.5})}{\log N} = \frac{\log(\varphi(m)/m) - 0.5\log N}{\log N} \to 0.5 \quad (N\to\infty).
\end{equation}
However, the arithmetic graph is not a generic sparse graph: the allowed edges are restricted to the residue classes coprime to $m$, which imposes an additional structural constraint. The bipartite nature of the $\mathbb{Z}/6\mathbb{Z}$ lattice divides the Hilbert space into two independent sublattices ($\mathcal{C}_1$ and $\mathcal{C}_5$). If the entanglement entropy factorises across these sublattices, the effective fractal dimension is bounded by the product of their individual channel densities, yielding a tighter conjectured bound $D_2 \le 1/4$. The empirical value $D_2 \approx 0.243$ for $m=6$ aligns closely with this conjecture, supporting the hypothesis that the arithmetic mask controls the fractal geometry down to the microscopic level.

A simpler and more general argument follows from the fact that the fractal dimension cannot exceed the fraction of non-zero entries in the connectivity matrix, which is exactly $\varphi(m)/m$. Thus $D_2 \le \varphi(m)/m$ for any $m$ and any $\nu\in(0.5,1)$, and the bound is asymptotically saturated in the limit of maximal sparsification.
\end{proof}

For $m=6$, $\varphi(6)/6 = 2/6 = 1/3$, and our numerical results give $D_2 \approx 0.243 < 1/3$, consistent with the bound. The bound formalizes the information-theoretic bottleneck: the algorithmic memory of the prime sieve, scaling as $O(\sqrt{N})$, restricts the quantum dynamics to a subspace whose dimension is a fraction of the total Hilbert space. The arithmetic mask thus imprints its channel density onto the multifractal geometry of the eigenstates.

\bibliographystyle{apsrev4-2}
\bibliography{references}

\end{document}